\documentclass[../../main.tex]{subfiles}

\graphicspath{{../../image/}{../../figure/}}

\begin{document}

% 原书物理页 15--30，印刷页码 XVII--XXXIII。
% 原书为无编号的卷首篇章，通篇为历史叙述，公式按原书在全篇内连续编号 (1)--(55)。
% 因处于章计数为 0 的卷首，align 自动给出与原书一致的编号形式 (1)、(2)……；
% 各编号公式均建 align: 标签，正文引用一律用 \eqref。
% 说明：本篇虽在叙述中顺带引入“绝热指数”“加速度波”“激波层”等术语，但其数学定义已在第 II、III 章
% 以正式定义给出，此处系历史叙述中的提及（v13 §3.0 E：历史叙述保持正文），故不再另立 definition 环境。

\chapter*{双曲守恒律早期历史概述}
\phantomsection\label{chapter:front-history}
\pdfbookmark[1]{双曲守恒律早期历史概述}{bm-front-history}
\markboth{双曲守恒律早期历史概述}{双曲守恒律早期历史概述}

双曲守恒律的一般理论乃至这个名称本身，不过是五十年前才出现的；然而这类偏微分方程组的特殊性质，早已在数学物理中出现的具体例子的语境里被辨认出来。此处的目的是追溯那些开创该领域并使之走上今日道路的早期奠基性工作。若干相关的经典论文已收入 Johnson 与 Chéret（1998）所编的文集。

接下来的叙述将描述：这门学科如何从流体动力学中产生，其早期步伐如何因热力学的混乱状态而受挫，如何被置于坚实的基础之上，以及它如何最终演化为偏微分方程理论的一个特殊分支。

本节可以独立于本书其余部分阅读，因为它本质上是自足的；但学生在熟悉正文所阐述的当前研究状况之后再回过头来读它，会获得额外的收益。据此，为凸显过去与现在的联系，此处的历史叙述得益于后见之明：自由地使用当代术语，取自原始文献的符号与方程已被转写为现代记号。

由于双曲守恒律的早期历史与气体动力学不可分割地交织在一起，我们先简要回顾十九世纪之交的理想气体理论。该主题的细节见 Truesdell（1980）的历史论著。

理想气体的状态由其密度 $\rho$、压力 $p$ 与（绝对）温度 $\theta$ 确定，它们由与 Boyle、Gay-Lussac 和 Mariotte 之名相联系的状态方程相互关联：
\begin{align}
p=R\rho\theta,\label{align:front-ideal-gas}
\end{align}
其中 $R$ 是普适气体常数。代替 $\rho$，同样可以使用其倒数 $u=1/\rho$，即比容。

定压比热或定容比热 $c_p$ 或 $c_u$，是在压力或比容保持不变而温度变化时，气体中所储存热量的变化率。比值 $\gamma=c_p/c_u$ 是一个大于 $1$ 的常数，称为\emph{绝热指数}。

正压热力学过程满足 $p=p(\rho)$，可在力学的范围内处理，而不必顾及温度。最简单的例子是等温过程：温度保持不变，于是由\eqref{align:front-ideal-gas}得
\begin{align}
p=a^{2}\rho.\label{align:front-isothermal}
\end{align}
更为微妙的是等熵过程或绝热过程\footnote{“绝热”一词由 Rankine 于 1859 年创造，他也是用符号 $\gamma$ 表示绝热指数之肇始者。不过在本书后文中我们将采用较新的术语“等熵”，而把“绝热”留给一个相关但不同的用法；见\cref{section:ch02-热弹性}。}：温度与比容同时按这样的比例变化，使得气体任何部分中所储存的热量保持不变。正如 Laplace 与 Poisson 所示，这一假设导致
\begin{align}
p=a^{2}\rho^{\gamma}.\label{align:front-isentropic}
\end{align}

双曲守恒律方程组最古老、至今仍最突出的范例，由正压气体流动的 Euler 方程组提供；这些方程表达质量与动量的守恒，把速度场 $v$ 与密度场、压力场联系起来。Euler 的相关论著——以其关于流体动力学的决定性表述为顶点——收入 Euler（1954），其中还载有 Truesdell 富有教益的评注。除以其名字命名的方程之外，Euler 还导出了如今称为无旋流动的 Bernoulli 方程；之所以这样称呼，是因为在定常流动中它归结为 Daniel Bernoulli 所发现的那个著名定律。我们将在本书正文中多次遇到上述方程，自\cref{subsection:ch03-等熵流}开始。

弹性流体中的内力由静水压力传递，而静水压力是标量场。因此，Euler 方程组构成一个具有独特几何结构的守恒律方程组。如\cref{chapter:ch01-平衡律}所讨论的、体现通量场张量本性的更一般类型的守恒律，于 1820 年代从 Cauchy 关于弹性理论的开创性工作中涌现（Cauchy，1823，1827，1829）。然而，如下文所示，关于双曲守恒律的早期工作几乎只处理一维空间的情形，而 Euler 方程组在这一情形下是一个完全具有代表性的例子。

在一篇关于声学理论的重要论文中，Poisson（1808）考察了理想气体直线等温流动的 Euler 方程组与 Bernoulli 方程，即
\begin{align}
\left\{\begin{aligned}
&\partial_t\rho+\partial_x(\rho v)=0\\
&\partial_t(\rho v)+\partial_x(\rho v^{2})+a^{2}\partial_x\rho=0,
\end{aligned}\right.\label{align:front-poisson-system}
\end{align}
\begin{align}
\partial_t\varphi+\frac{1}{2}(\partial_x\varphi)^{2}+a^{2}\log\rho=0,\label{align:front-bernoulli}
\end{align}
其中 $\varphi$ 是速度势，$v=\partial_x\varphi$。通过在\eqref{align:front-bernoulli}与\eqref{align:front-poisson-system}的第一式之间消去 $\rho$，他导出了仅含 $\varphi$ 的二阶方程
\begin{align}
\partial_t^{2}\varphi+2(\partial_x\varphi)(\partial_x\partial_t\varphi)+(\partial_x\varphi)^{2}\partial_x^{2}\varphi-a^{2}\partial_x^{2}\varphi=0.\label{align:front-second-order}
\end{align}

他运用 Laplace 与 Legendre 所创的一种解微分方程的方法，得出结论：任何满足函数方程
\begin{align}
\partial_x\varphi=f(x+(a-\partial_x\varphi)t),\label{align:front-functional}
\end{align}
（其中 $f$ 为任意光滑函数）的 $\varphi$，都是\eqref{align:front-second-order}的一个特解。

用当代术语来说，人们把 Poisson 的解识别为一条简单波（见\cref{section:ch07-简单波}），在其上 Riemann 不变量（见\cref{section:ch07-Riemann不变量}）$v+a\log\rho$ 为常数，于是 $v$ 满足方程
\begin{align}
\partial_t v+v\partial_x v-a\partial_x v=0,\label{align:front-simple-wave-eq}
\end{align}
该方程容许解
\begin{align}
v=f(x+(a-v)t),\label{align:front-simple-wave}
\end{align}
其中 $f$ 是任意光滑函数。

在 Poisson 的论文发表四十年后，英国天文学家 Challis（1848）作出如下观察：取 $f(x)=-\sin(\tfrac{1}{2}\pi x)$ 时，\eqref{align:front-simple-wave}给出沿直线 $x=-at$ 有 $v=0$，沿直线 $x=-1-(a-1)t$ 有 $v=1$；这引出一个悖论：在这两条直线的交点 $(-a,1)$ 处，$v$ 必须同时等于 $0$ 与 $1$。这是关于经典解破裂的最早记载，而经典解的破裂贯穿了整个双曲守恒律理论。

Challis 提出的问题几乎立即由他在 Cambridge 大学的同事 Stokes（1848）加以处理。Stokes 指出，按照 Poisson 的解\eqref{align:front-simple-wave}，沿每一条直线 $x=\bar{x}-(a-f(\bar{x}))t$ 有 $v(x,t)=f(\bar{x})$，并且
\begin{align}
\partial_x v(x,t)=\frac{f'(\bar{x})}{1+tf'(\bar{x})}.\label{align:front-stokes-derivative}
\end{align}
因此，除非 $f$ 非减，波将在 $t=-1/f'(\bar{x})$ 处破裂，其中 $f'(\bar{x})$ 是 $f'$ 的最小值。他随后思考奇性形成之后会发生什么，并提出一个原创而大胆的猜想。用他自己的话说：“试作最自然的设想，或许是形成一个间断曲面，通过它时密度与速度发生突变。”他似乎充分意识到这是一个意义深远的想法，远远超出 Poisson 解的特殊设定，因为他写道：“尽管我是通过考虑积分\eqref{align:front-simple-wave}的解释而被引向这一主题的，但此处引入对间断运动的考察，并非与那种解释相联系，而纯粹是为了它本身；我希望这两个主题被视为彼此完全独立。”

Stokes 接着刻画那些符合质量与动量守恒这两条支配物理定律的跳跃间断——在光滑流动的范围内，正是这些定律构成 Euler 方程组的基础。他假设密度与速度在一条以速度 $\sigma$（即斜率为 $\sigma$）传播的间断线两侧由 $(\rho_-,v_-)$ 跳跃到 $(\rho_+,v_+)$，并证明
\begin{align}
\left\{\begin{aligned}
&\rho_+v_+-\rho_-v_-=\sigma(\rho_+-\rho_-)\\
&\rho_+v_+^{2}+a^{2}\rho_+-\rho_-v_-^{2}-a^{2}\rho_-=\sigma(\rho_+v_+-\rho_-v_-).
\end{aligned}\right.\label{align:front-stokes-jump}
\end{align}
在上述两式之间消去 $\sigma$，他得到
\begin{align}
\rho_-\rho_+(v_+-v_-)^{2}=a^{2}(\rho_+-\rho_-)^{2}.\label{align:front-stokes-elimination}
\end{align}

于是，Stokes（1848）在等温流动的 Euler 方程组\eqref{align:front-poisson-system}的语境中引入了激波的概念，并导出了如今称为 Rankine--Hugoniot 跳跃条件的条件（见\cref{section:ch08-强激波弱激波中等强度激波}），后者刻画了\eqref{align:front-poisson-system}的分布弱解。这篇论文是双曲守恒律理论的基石之一。然而，这门学科的发展很快便撞上了路障。

Stokes 考虑带跳跃间断的流动的想法，先是受到 William Thomson 爵士（Lord Kelvin）——据称是私下地——后来则受到 Lord Rayleigh 的批评，后者既在私人通信中（Rayleigh，1877）也在出版物中（Rayleigh，1878，§253）提出批评，理由如下：他们论证说，跳跃间断不应产生或消耗机械能。一项计算表明，这将要求
\begin{align}
2\rho_-\rho_+\log\frac{\rho_-}{\rho_+}=\rho_-^{2}-\rho_+^{2},\label{align:front-rayleigh-condition}
\end{align}
而这与 $\rho_-\neq\rho_+$ 不相容。

为把上述论证置于守恒律理论的当下语境中，应当注意：Euler 方程组\eqref{align:front-poisson-system}的任一光滑解自动满足机械能守恒律
\begin{align}
\partial_t\left(\frac{1}{2}\rho v^{2}+a^{2}\rho\log\rho\right)+\partial_x\left(\frac{1}{2}\rho v^{3}+a^{2}\rho v\log\rho+a^{2}\rho v\right)=0.\label{align:front-energy-law}
\end{align}
用当代术语说，$\tfrac{1}{2}\rho v^{2}+a^{2}\rho\log\rho$ 是方程组\eqref{align:front-poisson-system}的一个熵，其熵流为 $\tfrac{1}{2}\rho v^{3}+a^{2}\rho v\log\rho+a^{2}\rho v$；见\cref{section:ch07-熵熵流对}。假设即使在激波存在时机械能也应守恒，就导出跳跃条件
\begin{align}
&\frac{1}{2}\rho_+v_+^{3}+a^{2}\rho_+v_+\log\rho_++a^{2}\rho_+v_+-\frac{1}{2}\rho_-v_-^{3}-a^{2}\rho_-v_-\log\rho_--a^{2}\rho_-v_-\label{align:front-energy-jump}\\
&\qquad=\sigma\left[\frac{1}{2}\rho_+v_+^{2}+a^{2}\rho_+\log\rho_+-\frac{1}{2}\rho_-v_-^{2}-a^{2}\rho_-\log\rho_-\right].\nonumber
\end{align}
在\eqref{align:front-stokes-jump}的第一式与\eqref{align:front-energy-jump}之间消去 $\sigma$，并利用\eqref{align:front-stokes-elimination}，便得到\eqref{align:front-rayleigh-condition}。

Stokes 被上述批评所说服，或许也受到了威慑。1877 年 6 月 5 日，他以一种歉意的语调回复 Lord Rayleigh，放弃了自己的理论：“感谢您指出您所提到的那篇论文中那种古怪运动的异议。W. Thomson 爵士许多年前就向我指出过同样的问题；若我有机会写下任何与该主题相关的东西，或者，若不借助那样的机会而我的论文竟引起了注意，我本应提到这一点。然而，为埋没在其他科学古董之中的一篇论文里的一段文字写批评，似乎几乎不值得”（Stokes，1966）。此外，当他的文集于 1880 年出版时，他删去了 1848 年论文中讨论间断流动的部分，并插入一则注解，大意是激波不可能存在，因为它们与\eqref{align:front-rayleigh-condition}所表达的能量守恒不相容（Stokes，1883）。关于 Stokes、Rayleigh 与 Thomson 之间通信的更详尽记述，连同对原始文献的指引，见 Salas（2007）。

Stokes 的平反不得不等待热力学的发展。然而令人费解的是，他竟如此轻易地准备放弃自己的理论，尽管如下文所示，到 1880 年该主题已取得实质性进展。

为按时间顺序追踪这些发展，让我们回到 1840 年代后期。皇家天文学家 Airy（1845）曾证明：浅水槽中纵向水波的传播，在 Lagrange 坐标下由二阶方程
\begin{align}
\partial_t^{2}w=a^{2}(1+\partial_x w)^{-3}\partial_x^{2}w\label{align:front-airy-long-wave}
\end{align}
所支配。通过令 $u=1+\partial_x w$、$v=\partial_t w$，可以把\eqref{align:front-airy-long-wave}改写为双曲守恒律方程组
\begin{align}
\left\{\begin{aligned}
&\partial_t u-\partial_x v=0\\
&\partial_t v+\partial_x\left(\frac{1}{2}a^{2}u^{-2}\right)=0.
\end{aligned}\right.\label{align:front-shallow-water}
\end{align}
这一方程组在 Euler 坐标下写出的形式称为浅水方程，它与绝热指数 $\gamma=2$ 的等熵流动的 Euler 方程组相同。其推导见\cref{section:ch07-一维平衡律}。

针对 Challis 1848 年的论文，Airy（1849）指出：类似地，等温流动的 Euler 方程组\eqref{align:front-poisson-system}在 Lagrange 坐标下可以改写为标量二阶方程
\begin{align}
\partial_t^{2}w=a^{2}(1+\partial_x w)^{-2}\partial_x^{2}w,\label{align:front-airy-isothermal}
\end{align}
其中 $w$ 是位移，于是 $u=1+\partial_x w$ 是比容，$v=\partial_t w$ 是速度。

Airy 随后指出，\eqref{align:front-airy-long-wave}与\eqref{align:front-airy-isothermal}分别容许解
\begin{align}
\partial_t w-2a(1+\partial_x w)^{-\frac{1}{2}}=\text{常数},\label{align:front-airy-invariant-1}
\end{align}
\begin{align}
\partial_t w-a\log(1+\partial_x w)=\text{常数},\label{align:front-airy-invariant-2}
\end{align}
其推导他归功于 De Morgan。这是 Riemann 不变量首次与简单波相联系而在出版物中明确出现。

参照 Stokes 的工作，Airy 猜想：长河中的涌潮以及广阔平坦沙滩上的碎浪，也许可以用他的浅水方程的解——带有跳跃间断的解——来表示。

下一个值得注意的贡献来自 Earnshaw（1860）。其分析的出发点是如下观察：若 $y(x,t)$ 对某个函数 $f$ 满足
\begin{align}
\partial_t y-f(\partial_x y)=\text{常数},\label{align:front-earnshaw-compatible}
\end{align}
则
\begin{align}
\partial_t^{2}y=[f'(\partial_x y)]^{2}\partial_x^{2}y.\label{align:front-earnshaw-second-order}
\end{align}
他把\eqref{align:front-earnshaw-second-order}与 Lagrange 坐标下流体直线正压运动的方程
\begin{align}
\rho_0\partial_t^{2}y+\partial_x p(\partial_x y)=0\label{align:front-earnshaw-motion}
\end{align}
相比较，其中 $y(x,t)$ 是流体粒子 $x$ 在时刻 $t$ 的位置，$\rho_0$ 是参考密度，$p$ 是压力。于是，可按\eqref{align:front-earnshaw-compatible}的形式得到\eqref{align:front-earnshaw-motion}的特解，其中 $f$ 由 $[f'(u)]^{2}=-\rho_0^{-1}p'(u)$ 算出。

取 $\rho_0=1$，按当代的惯例引入比容 $u=\partial_x y$ 与速度 $v=\partial_t y$，把\eqref{align:front-earnshaw-motion}改写为等价的一阶方程组
\begin{align}
\left\{\begin{aligned}
&\partial_t u-\partial_x v=0\\
&\partial_t v+\partial_x p(u)=0.
\end{aligned}\right.\label{align:front-earnshaw-system}
\end{align}
于是 $v-f(u)$ 是\eqref{align:front-earnshaw-system}的 Riemann 不变量，而解 $v-f(u)=\text{常数}$ 表示简单波。

Earnshaw 对等温流动（此时 $p=a^{2}\rho=a^{2}u^{-1}$）与等熵流动（此时 $p=a^{2}\rho^{\gamma}=a^{2}u^{-\gamma}$，$\gamma>1$）分别算出 $f(u)$，从而得到\eqref{align:front-earnshaw-compatible}形式的解。用当代记号，
\begin{align}
v\pm a\log u=\text{常数},\label{align:front-earnshaw-isothermal}
\end{align}
\begin{align}
v\pm\frac{2a\gamma^{\frac{1}{2}}}{\gamma-1}u^{\frac{\gamma-1}{2}}=\text{常数}.\label{align:front-earnshaw-isentropic}
\end{align}
依符号的不同，他把这些解设想为“凝聚”波或“稀疏”波。如今以“压缩”一词代替“凝聚”。

Earnshaw 的另一个有趣观察是：Euler 方程组容许行波解（用他的术语，“不发生类型改变的传播波”）当且仅当 $p=a-b/\rho$，这辨识出如今所称的 Chaplygin 气体；见\cref{section:ch02-热弹性}。这是与线性退化的非线性系统的最早相遇；见\cref{section:ch07-真正非线性}。Earnshaw 进一步指出，对于这个特殊的 $p$，Lagrange 坐标下的运动方程\eqref{align:front-earnshaw-motion}是线性的。

现在我们来到双曲守恒律理论发展中的一个里程碑，即 Riemann（1860）那篇著名的论文。这是最早从分析的角度处理 Euler 方程组的工作。然而应当指出，作者不遗余力地从物理学出发为他的工作提供动机。他表示希望：他的结果除纯数学意义之外，还能在声学中找到应用。

Riemann 考虑状态方程 $p=p(\rho)$ 一般的气体直线正压流动，只要求 $p'(\rho)>0$，并把 Euler 方程组写成
\begin{align}
\left\{\begin{aligned}
&\partial_t\log\rho+v\partial_x\log\rho=-\partial_x v\\
&\partial_t v+v\partial_x v=-p'(\rho)\partial_x\log\rho.
\end{aligned}\right.\label{align:front-riemann-system}
\end{align}
如上所见，如今所称的 Riemann 不变量已出现在 Poisson、Airy 与 Earnshaw 的工作中，只是完全与简单波相联系。与此相对，Riemann 在\eqref{align:front-riemann-system}的任一光滑解的语境中定义他的不变量
\begin{align}
r=\frac{1}{2}\left[v+\int\sqrt{p'(\rho)}\,\mathrm{d}\rho\right],\qquad
s=\frac{1}{2}\left[-v+\int\sqrt{p'(\rho)}\,\mathrm{d}\rho\right],\label{align:front-riemann-invariants}
\end{align}
并证明它们满足
\begin{align}
\partial_t r=-\left(v+\sqrt{p'(\rho)}\right)\partial_x r,\qquad
\partial_t s=-\left(v-\sqrt{p'(\rho)}\right)\partial_x s,\label{align:front-riemann-transport}
\end{align}
于是 $r$ 与 $s$ 分别沿如今所称的前向特征线与后向特征线保持常数。

Riemann 随后发明了如今所称的速矢变换（hodograph transformation，见\cref{section:ch12-熵熵流对与速矢变换}），它通过颠倒 $(x,t)$ 与 $(r,s)$ 作为自变量与因变量的角色，把非线性方程组\eqref{align:front-riemann-transport}改写为线性方程。为此，在注意到
\begin{align}
\partial_r\left[x-\left(v-\sqrt{p'(\rho)}\right)t\right]=-\partial_s\left[x-\left(v+\sqrt{p'(\rho)}\right)t\right]\label{align:front-hodograph-identity}
\end{align}
之后，他引入势 $w(r,s)$，
\begin{align}
\partial_r w=x-\left(v+\sqrt{p'(\rho)}\right)t,\qquad
-\partial_s w=x-\left(v-\sqrt{p'(\rho)}\right)t,\label{align:front-hodograph-potential}
\end{align}
并证明它满足线性方程
\begin{align}
\partial_r\partial_s w=m(\partial_r w+\partial_s w),\label{align:front-hodograph-eq}
\end{align}
其中 $m$ 是由 $p(\rho)$ 诱导出的 $\rho$ 的某个函数，从而只依赖于 $r+s$。特别地，在等温情形 $p=a^{2}\rho$ 下，$m=-(2a)^{-1}=$ 常数。人们须通过解\eqref{align:front-hodograph-eq}确定 $w(r,s)$，再由\eqref{align:front-hodograph-potential}得到 $x(r,s)$ 与 $t(r,s)$，最后反演这些函数以得到 $r(x,t)$ 与 $s(x,t)$。

接着，Riemann 在附加假设 $p''(\rho)\geqslant0$ 下考虑带跳跃间断的解。通过在激波 $x=\xi(t)$（左状态 $(\rho_-,v_-)$、右状态 $(\rho_+,v_+)$、速度 $\dot{\xi}$）两侧平衡质量与动量，他导出了如下形式的跳跃条件
\begin{align}
\dot{\xi}=v_-\pm\left[\frac{\rho_+}{\rho_-}\,\frac{p(\rho_+)-p(\rho_-)}{\rho_+-\rho_-}\right]^{\frac{1}{2}}
=v_+\pm\left[\frac{\rho_-}{\rho_+}\,\frac{p(\rho_+)-p(\rho_-)}{\rho_+-\rho_-}\right]^{\frac{1}{2}}.\label{align:front-riemann-jump}
\end{align}
他没有论及能量守恒问题，因为他大概并不知道 Kelvin 与 Rayleigh 对 Stokes 工作所提出的异议。尽管如此，他强调：基于稳定性的理由，只有压缩激波在物理上有意义。事实上，他以如今所称的 Lax $E$-条件（见\cref{section:ch08-lax容许性准则}）的形式假定激波可容许性，即对前向激波有
\begin{align}
v_++\sqrt{p'(\rho_+)}<\dot{\xi}<v_-+\sqrt{p'(\rho_-)},\label{align:front-lax-e}
\end{align}
对后向激波类似。他还构造了如今所称的激波曲线或 Hugoniot 轨迹（见\cref{section:ch08-hugoniot轨迹}），并描述其形状。

最后，Riemann 引入并求解了如今以他命名的那个著名问题（见\cref{chapter:ch09-容许波扇与riemann问题}），即他证明：状态变量 $(\rho,v)$ 的跳跃一般被分解为一个由后向波与前向波组成的出射波扇，其中每一个波或为压缩激波，或为中心稀疏简单波。

Riemann 这篇非凡的论文为一维空间中双曲守恒律方程组的一般理论奠定了基础。如下文所示，在同时守恒能量、质量与动量的气体流动中，激波不可能是等温的或等熵的。这引出了文献中广泛流传的一种抱怨，即 Riemann 对激波的处理是有缺陷的。我们已经看到，正是同一个论证迫使 Stokes 放弃了他的等温激波理论。然而，这样的批评是不得要领的。Riemann 关于跳跃条件的方程\eqref{align:front-riemann-jump}在等熵气体动力学的框架内是数学上准确的，而后者是对更完备的热力学理论的一个在物理上正当的简化。

1860 年之后，激波的概念逐渐获得接受。到 1870 年，从 Rankine（1870）的一则脚注来看，Kelvin——Stokes 想法的早期批评者之一——已准备承认压缩激波，同时以不稳定为由拒绝稀疏激波。

1877 年，Christoffel（1877）考虑了三维空间中正压气体流动的激波，并导出了相应的跳跃条件。事实上，用方括号括起场的符号以表示其跨激波的跳跃这一通行做法，正是由他引入的。见 Hölder（1981）。

1875 年至 1889 年间，Ernst Mach 在其学生的协助下，在布拉格的德国大学完成了最早的激波实验。他利用由 Leiden 瓶放电触发的电火花，在两块平行玻璃板之间所夹的空气层中产生激波，其中一块玻璃板上覆有烟炱。随后根据激波在烟炱上留下的痕迹确定其路径。借助这一装置，他成功测量了激波的速度，表明它们是超声速的。他还观察到，激波对刚性边界的斜入射可以诱发正规反射或不正规反射（现称 Mach 反射）。Mach 反射现象也在液体的表面波中被检测到（水银、牛奶、糖浆）；激波是通过把一段 V 形铁丝投落到液面上产生的。后来，借助威力更强的光学纹影法，Mach 及其学生证明：以超声速穿过空气的弹体周围环绕着弓形激波。关于 Mach 实验的历史记述，连同对原始论文的指引——这些论文发表于维也纳科学院 Sitzungsherichte 第 72--92 卷——见 Krehl 与 van der Geest（1991）。

十九世纪中叶是掌握热力学原理取得迅速进展的时期。内能 $\varepsilon$ 的核心作用得到认识，能量守恒定律得以确立。熵的概念由 Clausius（1854）引入并命名，热力学第二定律也得到表述。不久，比熵 $s$ 便在热力学变量的清单上占据了位置，与 $\varepsilon$、$\theta$、$p$、$u$ 和 $\rho$ 并列。一般的热状态方程 $\varepsilon=\varepsilon(\rho,\theta)$ 由 Kirchhoff（1868）考察，而等价的量热状态方程 $\varepsilon=\varepsilon(u,s)$ 由 Gibbs（1873）引入，他假定了如今以他命名的规则：
\begin{align}
\mathrm{d}\varepsilon=\theta\,\mathrm{d}s-p\,\mathrm{d}u.\label{align:front-gibbs}
\end{align}
然而，如下文所示，大多数作者仍坚持理想气体这一特殊情形，其状态方程为
\begin{align}
\varepsilon=\frac{R}{\gamma-1}\theta,\qquad p=R\rho\theta,\qquad
s=\frac{R}{\gamma-1}\log\frac{\theta}{\rho^{\gamma-1}}.\label{align:front-ideal-gas-state}
\end{align}

就流体动力学而言，人们认识到 Euler 方程组必须补充一个独立的场方程，以表达（机械能与热能合并的）能量守恒。这一方程由 Kirchhoff（1868，1894）对热粘弹性、导热气体导出。在无黏性与热传导时，把能量守恒方程附加到 Euler 方程组上，便得到双曲守恒律方程组的另一个重要范例；本书正文将多次遇到它，自\cref{subsection:ch03-热弹性流体}开始。

确定在激波存在时表达能量守恒的跳跃条件这一任务，由 Rankine（1870）与 Hugoniot（1887）承担。Rankine 针对向未受扰介质中运动的激波导出其跳跃条件：把激波处机械能的损失与热流相平衡，使总能量生成量为零。要跟上他的讨论并不容易\footnote{关于 Rankine 论证的重构，见 Rayleigh（1910）。}——它基于物理论证——但最终他得到了正确的方程。

与 Rankine 的途径相反，Hugoniot（1887）的处理具有清晰的数学风格。他在 Lagrange 坐标下考虑一般非线性弹性介质的直线运动。首先，他详细讨论了正压运动，它由标量二阶方程\eqref{align:front-earnshaw-motion}支配，或等价地由一阶方程组\eqref{align:front-earnshaw-system}支配。他重述（抑或重新发明？）了 Earnshaw 与 Riemann 的结果，却未引用这两位作者。特别地，他按与 Earnshaw 相同的步骤确定了两族 Riemann 不变量，并随后像 Riemann 早已做过的那样证明它们沿同族特征线为常数。他还指出了 Riemann 不变量与简单波之间的联系。

Hugoniot 具有开拓性的贡献见于其论文的第五章，那里他讨论激波，尤其是 §§149--158，其中他导出了由能量守恒所规定的跳跃条件。正是在那里，人们第一次遇到 Lagrange 坐标下完整的一组跳跃条件：
\begin{align}
\left\{\begin{aligned}
&\sigma(u_+-u_-)=v_--v_+\\
&\sigma(v_+-v_-)=p_+-p_-\\
&\sigma\left(\varepsilon_++\frac{1}{2}v_+^{2}-\varepsilon_--\frac{1}{2}v_-^{2}\right)=p_+v_+-p_-v_-,
\end{aligned}\right.\label{align:front-hugoniot-system}
\end{align}
这是对理想气体写出的，其内能 $\varepsilon=pu/(\gamma-1)$；见\eqref{align:front-ideal-gas-state}。把\eqref{align:front-hugoniot-system}中的三个方程结合起来，便得到著名的跳跃条件
\begin{align}
\varepsilon_+-\varepsilon_-+\frac{1}{2}(p_++p_-)(u_+-u_-)=0,\label{align:front-hugoniot-condition}
\end{align}
它不含 $v$ 或 $\sigma$。Hugoniot 对理想气体的情形以如下形式导出这一方程：
\begin{align}
\frac{p_+}{p_-}=\frac{2u_-+(\gamma-1)(u_--u_+)}{2u_++(\gamma-1)(u_+-u_-)}.\label{align:front-hugoniot-isentrope}
\end{align}
他似乎并不知道 Rankine（1870）已得到过类似的结果。

上述方程在气体动力学理论及其应用中产生了巨大影响。因此，如今所有与激波相联系的跳跃条件被统称为“Rankine--Hugoniot 跳跃条件”，尽管“Stokes 跳跃条件”这一名称会更准确地反映历史记录。

Hugoniot（1885）的另一项重要贡献是引入并研究弱波（又称加速度波），即传播的特征曲面，其两侧状态变量本身连续而其导数经历跳跃间断。弱波早已出现在 Euler 的声学研究之中。

到二十世纪之交，双曲守恒律领域正在向外分枝，在燃烧、爆轰与空气动力学的科学与技术中找到应用。我们将不追索那些新方向上的发展史，因为此处的目的是追溯导致数学理论今日状况的核心思想的演进。关于爆轰这一主题在二十世纪初的状况，见 Jouguet（1917）的著作。要熟悉空气动力学浩繁的文献，一个良好的起点是 von Mises（1958）的教材，其中载有历史性的参考文献与一份庞大的书目。

1900 年前后双曲守恒律基本理论的研究状况，可由 Duhem（1903）、Hadamard（1903）与 Weber（1901）的著作来例示。Weber 主要对上述 Riemann 的论文作了详细阐发。就 Hadamard 而言，他把 Riemann 的工作与 Hugoniot 的结果结合起来加以陈述。Duhem 与 Hadamard 都对多空间维中激波与弱波的传播作了详尽阐述。特别地，Hadamard 作出了一个重要观察：无旋流动穿过弱波后仍保持无旋，而穿过非平面激波后则获得涡量。Duhem 强调热力学的蕴涵。他以如今以他命名的著名场不等式的形式假定第二定律（见\cref{section:ch02-连续介质热力学平衡律}）：
\begin{align}
\rho_0\dot{s}+\operatorname{Div}\left(\frac{1}{\theta}Q\right)\geqslant0,\label{align:front-duhem-inequality}
\end{align}
其中 $Q$ 是热流向量。此外，他证明：在无黏性与热传导时，质量、动量与能量的守恒律方程组连同 Gibbs 规则\eqref{align:front-gibbs}，蕴涵光滑的热力学过程必然等熵：$\dot{s}=0$。

随着新世纪的到来，理论面临着激波的物理可容许性问题。读者或许记得，在上个世纪，自 Riemann 起若干作者已赞同这样的观点：压缩激波是稳定的，从而是可容许的；与此相反，稀疏激波显然不稳定，因为它们易于瓦解为稀疏简单波。然而，激波稳定性与热力学第二定律之间的联系仍难以捉摸。另一个争论之点是正压流动中激波的物理地位，鉴于 Rayleigh 的论证：它们不能守恒机械能。

处理这些问题的第一步由 Jouguet（1901，1904）迈出，他们证明：在理想气体中，只有压缩激波与热力学第二定律相容。确实，在无热流 $Q=0$ 时，\eqref{align:front-duhem-inequality}归结为 $\dot{s}\geqslant0$，而以速度 $\sigma$ 传播的激波所产生的 $\dot{s}$ 与 $\sigma(s_--s_+)$ 同号。借助\eqref{align:front-ideal-gas-state}，
\begin{align}
s_+-s_-=\frac{R}{\gamma-1}\log\left(\frac{p_+u_+^{\gamma}}{p_-u_-^{\gamma}}\right).\label{align:front-entropy-jump}
\end{align}
此外，$p_+/p_-$ 与 $u_+/u_-$ 由 Hugoniot 的方程\eqref{align:front-hugoniot-isentrope}相联系。于是容易看出，$s_+-s_-$ 与 $u_+-u_-$ 同号，从而与 $\rho_+-\rho_-$ 反号。

在上述工作之后，Duhem（1909）证明：在状态方程为 $p=p(u,s)$ 的任何气体中，弱压缩激波事实上满足热力学第二定律，只需约束 $p_{uu}>0$。在此过程中，他作出了一个重要观察：跨弱激波的熵跃变在激波强度上是三阶小量。他进一步证明：弱压缩激波相对于较密的流体是亚声速的，相对于较稀的流体是超声速的。因此，用当代术语说，Duhem 证明了：当气体动力学的守恒律方程组是真正非线性时，Lax $E$-条件是弱激波的热力学第二定律的体现；见\cref{section:ch08-熵激波容许性准则}。

与 Jouguet 的工作平行，Zemplén（1905）研究了理想气体中激波所产生的机械能与热能之间的平衡，并证明：在压缩激波处机械能转化为热，而在稀疏激波处这一过程将会反向。因此，热力学第二定律只允许把机械能转化为热的激波。

关于正压流动 $p=p(\rho)$，Burtton（1893）、Weber（1901）与 Rayleigh（1908）最终都准备在激波服从热力学第二定律的前提下赋予其物理地位，而他们把该定律解释为要求机械能的生成非正。Weber 证明：当 $p''(\rho)\geqslant0$ 时，只有压缩激波满足这一要求。他在 Euler 坐标下工作，但为避免写出过多新方程，此处我们把他计算转写到 Lagrange 坐标：以速度 $\sigma$ 传播的激波产生机械能的速率是
\begin{align}
\dot{E}=-\sigma\left\{\frac{1}{2}v_+^{2}-\frac{1}{2}v_-^{2}-\int_{u_-}^{u_+}p(u)\,\mathrm{d}u\right\}+p_+v_+-p_-v_-.\label{align:front-weber-production}
\end{align}
由于质量与动量守恒，$u_\pm$、$p_\pm$ 与 $v_\pm$ 由\eqref{align:front-hugoniot-system}中的前两个跳跃条件相联系；借助这些关系，\eqref{align:front-weber-production}可写成
\begin{align}
\dot{E}=\sigma\left\{\int_{u_-}^{u_+}p(u)\,\mathrm{d}u-\frac{1}{2}(p_++p_-)(u_+-u_-)\right\}.\label{align:front-weber-energy}
\end{align}
现在清楚了：假设 $p''(u)>0$，则 $\dot{E}<0$ 当且仅当 $\sigma(u_+-u_-)>0$，即 $\sigma(\rho_+-\rho_-)<0$。

用当代术语说，Weber 是把机械能 $\tfrac{1}{2}v^{2}-\int p(u)\,\mathrm{d}u$ 用作方程组\eqref{align:front-earnshaw-system}的凸熵，相应的熵流为 $vp(u)$，并证明：当方程组真正非线性时，熵可容许性条件等价于 Lax $E$-条件；见\cref{section:ch08-熵激波容许性准则}。

辨识物理上可容许激波的另一种（尽管相关的）途径是“消失黏性”方法。Stokes（1848）与 Hugoniot（1887）等人已意识到，黏性和/或热传导会把激波抹平。这一点由 Duhem（1903）加以形式化。于是，激波处所发生的机械能损失可归因于黏性与热传导所诱发的“内摩擦”的作用。取零黏性极限的合理性，可以通过证明物理上可容许的激波能与具有相同端状态和相同速度的黏性行波相配对来确立。通过改变坐标，人们可以只考虑定常激波与定常黏性波。

导热黏性流体直线流动的质量、动量与能量守恒律为
\begin{align}
\left\{\begin{aligned}
&\partial_t\rho+\partial_x(\rho v)=0\\
&\partial_t(\rho v)+\partial_x(\rho v^{2}+p-\mu\partial_x v)=0\\
&\partial_t\left(\rho\varepsilon+\frac{1}{2}\rho v^{2}\right)
+\partial_x\left(\rho v\varepsilon+\frac{1}{2}\rho v^{3}+pv-\mu v\partial_x v-k\partial_x\theta\right)=0,
\end{aligned}\right.\label{align:front-viscous-system}
\end{align}
其中 $\mu$ 是黏性系数，$k$ 是导热系数。在无黏性与热传导时，定常激波满足跳跃条件
\begin{align}
\left\{\begin{aligned}
&\rho_+v_+=\rho_-v_-=m\\
&mv_++p_+=mv_-+p_-=a\\
&m\varepsilon_++\frac{1}{2}mv_+^{2}+p_+v_+=m\varepsilon_-+\frac{1}{2}mv_-^{2}+p_-v_-=b.
\end{aligned}\right.\label{align:front-stationary-jump}
\end{align}
目标是确定激波层（也称激波剖面或激波结构），即\eqref{align:front-viscous-system}在 $(-\infty,\infty)$ 上、对给定参数 $m$、$a$ 与 $b$ 的定常解
\begin{align}
\left\{\begin{aligned}
&\rho v=m\\
&\rho v^{2}+p-\mu v'=a\\
&\rho v\varepsilon+\frac{1}{2}\rho v^{3}+pv-\mu vv'-k\theta'=b
\end{aligned}\right.\label{align:front-shock-profile}
\end{align}
此处及以下，撇号表示关于 $x$ 的微分。

上述计划几乎同时地、且看来彼此独立地由Rayleigh（1910）\footnote{除处理当前问题之外，Rayleigh 的论文还提供了对十九世纪激波理论发展的有趣评述。} 与 G. I. G. I. Taylor（1910）启动。他们的途径惊人地相似。假设气体是理想的（\eqref{align:front-ideal-gas-state}），他们在\eqref{align:front-shock-profile}的各方程之间消去 $p$、$\rho$、$\varepsilon$ 与 $\theta$，最终得到仅含 $v$ 的二阶方程：
\begin{align}
\frac{k\mu}{Rm}(vv')'=\left(\frac{2k}{R}+\frac{\mu}{\gamma-1}\right)vv'-\frac{ka}{Rm}v'
-\frac{\gamma+1}{2(\gamma-1)}mv^{2}+\frac{a\gamma}{\gamma-1}v-b.\label{align:front-taylor-equation}
\end{align}
当 $k=0$ 或 $\mu=0$ 时，Taylor 以闭形式解出\eqref{align:front-taylor-equation}。他还导出了 $k$ 与 $\mu$ 均为正但 $|v_+-v_-|$ 很小时解的渐近形式。特别地，他指出只有压缩激波才可能有黏性剖面。

Becker（1922）处理了同一问题，仍针对理想气体：他注意到，由\eqref{align:front-shock-profile}可知，温度作为比容的函数满足一个一阶微分方程，而在 $k/\mu$ 等于定压比热 $c_p$ 的特殊情形下，该方程可以闭形式积分。

第二次世界大战期间，来自各个专业领域的著名物理学家与数学家被吸引到气体动力学。当时的一个重要问题是真实气体在高温高压下、超出多方模型适用范围时的行为。与此相关，Bethe（1942）证明：在状态方程为 $p=p(u,s)$ 的气体中，仅凭条件 $p_{uu}>0$ 并不能保证任意强度的压缩激波与热力学第二定律相容且稳定。为此，必须附加假设，例如 $up_s+2\theta>0$ 与 $\theta p_u+pp_u<0$，而这类假设在气体发生相变时可能失效。

同一时期的另一项重要贡献是 Weyl（1949）的工作，它把由 Rayleigh（1910）、G. I. G. I. Taylor（1910）与 Becker（1922）所开创的激波层研究推广到具有一般状态方程的气体。通过组合\eqref{align:front-shock-profile}中的各方程，Weyl 在变量 $(u,\theta)$ 下导出如下一阶方程组
\begin{align}
\left\{\begin{aligned}
&\mu mu'=m^{2}u+p-a\\
&k\theta'=m\varepsilon-\frac{1}{2}m^{3}u^{2}+amu-b,
\end{aligned}\right.\label{align:front-weyl-system}
\end{align}
其中 $u$ 是比容 $\rho^{-1}$。注意到由\eqref{align:front-stationary-jump}，$(u_\pm,\theta_\pm)$ 是\eqref{align:front-weyl-system}的平衡点，他把激波剖面实现为连接 $(u_-,\theta_-)$（一个不稳定结点）与 $(u_+,\theta_+)$（一个鞍点）的轨道。两年后 Gilbarg（1951）给出了这一问题的决定性处理，对任意正的 $\mu$、$k$ 建立了激波层的存在性，并证明它收敛于：(a) 当 $\mu$ 与 $k$ 同时趋于零时的一个激波；(b) 当 $\mu$ 趋于零而 $k$ 保持固定时的一个连续激波层；(c) 当 $k$ 趋于零而 $\mu$ 保持固定时的一个一般说来不连续的激波层。正如我们将在\cref{section:ch08-黏性激波剖面}中看到的，Weyl 的途径如今已成为双曲守恒律一般理论中的标准做法。

作为战时工作的一部分而对该领域作出贡献的著名科学家中，包括 von Neumann。他就气体动力学中的激波理论准备了若干综述性报告 von Neumann（1963）、von Neumann（1963）、von Neumann（1963），其中载有许多富有洞见的观察。特别地，他详细阐述了斜激波反射问题，重新激活并普及了 Mach 在十九世纪的贡献。他还倡导借助科学计算获得解的思想；见 von Neumann（1963）。1949 年 8 月 17 日由 von Neumann（1963）主持、有 Burgers、Heisenberg、von Karman 等专家参加的一次专题讨论会的会议记录，为我们提供了当时被视为重大未决问题之一瞥。值得注意的是，那里提出的许多问题至今仍未解决。

到 1940 年代后期，关于双曲守恒律的大量信息已被积累起来，主要以气体动力学的形式出现。这些信息由数学家、物理学家、化学家与工程师在 150 年间导出，并以各式各样的风格与严密程度呈现。整合这些材料的任务由 Courant 与 Friedrichs（1948）承担，他们以数学语言对该主题作了权威性的综合。他们的著作曾经——并继续——在向数学界传播该理论的物理基础方面发挥重要作用。

1950 年之后，顺应偏微分方程这一大领域的当代趋势，双曲守恒律的研究聚焦于 Cauchy 问题的定性理论。建立经典解局部存在性的种子，已在 1930 年代由 Schauder（1935）关于二阶拟线性双曲方程的工作播下。Schauder 的策略——对线性化方程解的导数运用一系列 $L^{2}$“能量”估计，然后经由不动点论证过渡到拟线性情形——在过去数十年间被广泛利用，最终形成了可对称化守恒律方程组 Cauchy 问题的完备理论，即本书\cref{chapter:ch05-熵与经典解稳定性}所阐述的内容。

由于 Cauchy 问题的经典解一般在有限时间内破裂，人们至多只能期望建立大范围弱解的存在性。这一方向上的首次成功尝试见于 Hopf（1950）那篇奠基性的论文，它处理最简单的非线性标量守恒律
\begin{align}
\partial_t u+\partial_x\left(\frac{1}{2}u^{2}\right)=0\label{align:front-burgers}
\end{align}
的 Cauchy 问题。这一方程最初由 Bateman（1915）在一份不起眼的出版物中提出，作为气体动力学守恒律方程组的一个简单模型。它独立地重新出现在 Burgers（1940）关于湍流的工作中，如今被普遍称为 Burgers 方程。

采纳 Burgers 的观点，Hopf 把\eqref{align:front-burgers}视为带黏性的 Burgers 方程
\begin{align}
\partial_t u+\partial_x\left(\frac{1}{2}u^{2}\right)=\mu\partial_x^{2}u\label{align:front-burgers-viscous}
\end{align}
在 $\mu\downarrow0$ 时的极限。通过运用著名的 Hopf--Cole 变换
\begin{align}
u=-2\mu\partial_x\log\varphi,\label{align:front-hopf-cole}
\end{align}
——之所以这样称呼，是因为它也被 Cole（1951）独立发现——他把\eqref{align:front-burgers-viscous}归结为经典的热方程
\begin{align}
\partial_t\varphi=\mu\partial_x^{2}\varphi.\label{align:front-heat}
\end{align}
这使他能以显式形式解出方程\eqref{align:front-burgers-viscous}带初值 $u_0$ 的 Cauchy 问题：
\begin{align}
u(x,t)=\frac{\displaystyle\int_{-\infty}^{\infty}\frac{x-y}{t}\exp\left\{-\frac{1}{2\mu}F(x,y,t)\right\}\mathrm{d}y}
{\displaystyle\int_{-\infty}^{\infty}\exp\left\{-\frac{1}{2\mu}F(x,y,t)\right\}\mathrm{d}y},\label{align:front-hopf-solution}
\end{align}
其中
\begin{align}
F(x,y,t)=\frac{(x-y)^{2}}{2t}+\int_{0}^{y}u_0(z)\,\mathrm{d}z.\label{align:front-hopf-f}
\end{align}

在\eqref{align:front-hopf-solution}中令 $\mu\downarrow0$，他得到 Burgers 方程\eqref{align:front-burgers}的 Cauchy 问题的一个弱解 $u(x,t)$，它由初值 $u_0(x)$ 通过
\begin{align}
u(x,t)=\frac{x-y(x,t)}{t},\qquad\text{a.e. 于 }(-\infty,\infty)\times(0,\infty),\label{align:front-hopf-formula}
\end{align}
显式确定，其中 $y(x,t)$ 是函数 $F(x,y,t)$ 关于 $y$ 在 $(-\infty,\infty)$ 上的极小点。他还考察了这一解的几何结构与长时间行为。

Hopf 的论文激发了对标量守恒律的密集研究，起初在一维、最终在多空间维，从而产生了\cref{chapter:ch06-标量守恒律L1理论}与\cref{chapter:ch11-一维标量守恒律}所呈现的丰富理论。

下一个重大里程碑——它标志着这一历史性引言的结束——是 Lax（1957）那篇划时代的论文，它创造了“双曲守恒律”这一术语，并使该领域作为偏微分方程理论中一个新的主要分支而诞生。这是通过提炼、推广并形式化上文所述多年来在特殊方程组的语境中积累起来的原始材料而实现的。

Lax 论文的第一部分推广了上述 Hopf 的工作，特别地设计出\eqref{align:front-hopf-formula}的一个推广，用以求解一般凸标量守恒律的 Cauchy 问题。读者可在\cref{section:ch11-分界线不变量与Lax公式}中找到关于这一理论的叙述。

论文的第二部分通过引入严格双曲性、真正非线性、Riemann 不变量、简单波与 Lax $E$-条件等概念，为一维空间中双曲守恒律方程组的一般理论奠定了基础；这些概念在激波曲线与稀疏波曲线的构造以及 Riemann 问题的求解中汇合在一起。在很大程度上，正如本书\cref{chapter:ch07-一维双曲平衡律组}、\cref{chapter:ch08-容许激波}与\cref{chapter:ch09-容许波扇与riemann问题}所呈现的，一维双曲守恒律方程组理论的当今状况正是上述主题的展开。可以说，Lax 的论文为过去五十年间双曲守恒律领域的发展指明了方向。

\end{document}